\documentclass[11pt,a4paper]{article}
\usepackage{fontspec}
\IfFontExistsTF{texgyrepagella-regular.otf}
 {\setmainfont{texgyrepagella}[Extension=.otf,UprightFont=*-regular,BoldFont=*-bold,
   ItalicFont=*-italic,BoldItalicFont=*-bolditalic]}
 {\IfFontExistsTF{Palatino Linotype}{\setmainfont{Palatino Linotype}}{\setmainfont{Latin Modern Roman}}}
\linespread{1.05}
\usepackage[english]{babel}
\usepackage{amsmath,amssymb,amsthm,mathtools}
\usepackage{unicode-math}
\IfFontExistsTF{texgyrepagella-math.otf}%
{\setmathfont{texgyrepagella-math.otf}}{\setmathfont{Latin Modern Math}}
\usepackage[margin=26mm]{geometry}
\IfFileExists{microtype.sty}{\usepackage{microtype}}{}
\IfFileExists{titlesec.sty}{\usepackage[nobottomtitles*]{titlesec}
  \renewcommand{\bottomtitlespace}{.12\textheight}}{}
\usepackage{booktabs,array}
\usepackage[shortlabels]{enumitem}
\usepackage{tikz}
\usetikzlibrary{arrows.meta,positioning}
\usepackage{xcolor}
\usepackage[colorlinks=true,linkcolor={blue!45!black},citecolor={green!35!black},%
urlcolor={blue!55!black}]{hyperref}

\AtBeginDocument{\renewcommand{\tablename}{Table}}
\setlength{\parindent}{1.15em}
\setlength{\parskip}{0pt plus .35pt}
\setlength{\textfloatsep}{14pt plus 3pt minus 3pt}
\setlength{\intextsep}{12pt plus 3pt minus 3pt}
\setlength{\abovecaptionskip}{6pt}
\setlength{\belowcaptionskip}{2pt}
\makeatletter % float barrier, as \FloatBarrier in placeins
\providecommand\FloatBarrier{\par\begingroup\let\@elt\relax
  \edef\@tempa{\@botlist\@deferlist\@dbldeferlist}%
  \ifx\@tempa\@empty\endgroup\else\endgroup\newpage\fi\suppressfloats[t]}
\makeatother
\AtBeginDocument{\setlength{\abovedisplayskip}{9pt plus 2pt minus 3pt}%
  \setlength{\belowdisplayskip}{9pt plus 2pt minus 3pt}}
\clubpenalty=7000
\widowpenalty=7000
\displaywidowpenalty=3000
\definecolor{btinta}{HTML}{223B49}
\definecolor{bgris}{HTML}{DDE6E7}
\definecolor{bpapel}{HTML}{FFFFFF}
\newcommand{\bcuadricula}[1]{%
  \fill[bpapel] (0,0) rectangle (#1,#1);
  \foreach \r in {1,...,#1}{%
    \foreach \c in {1,...,#1}{%
      \ifodd\numexpr\r+\c\relax
        \fill[bgris] (\c-1,#1-\r) rectangle (\c,#1-\r+1);
      \fi
    }
    \node[font=\scriptsize,anchor=east] at (-.10,#1-\r+.5) {\r};
    \node[font=\scriptsize,anchor=south] at (\r-.5,#1+.06) {\r};
  }
  \foreach \t in {0,...,#1}{\draw[black!45,line width=.3pt] (\t,0)--(\t,#1) (0,\t)--(#1,\t);}
  \draw[btinta,line width=.6pt] (0,0) rectangle (#1,#1);
}
\newcommand{\breina}[3]{%
  \begin{scope}[shift={(#3-.5,#1-#2+.50)}]
    \fill[bpapel] (0,0) circle[radius=.36];
    \fill[btinta] (-.29,.15)--(-.13,.04)--(0,.30)--(.13,.04)--(.29,.15)--(.18,-.15)--(-.18,-.15)--cycle;
    \fill[btinta] (-.21,-.22) rectangle (.21,-.18);
  \end{scope}
}
\newcommand{\bmarca}[3]{\fill[btinta!62] (#3-.5,#1-#2+.50) circle[radius=.055];}

\setlength{\emergencystretch}{2em}
\allowdisplaybreaks[2]
\numberwithin{equation}{section}

\theoremstyle{plain}
\newtheorem{teorema}{Theorem}[section]
\newtheorem{proposicion}[teorema]{Proposition}
\newtheorem{lema}[teorema]{Lemma}
\newtheorem{corolario}[teorema]{Corollary}
\theoremstyle{definition}
\newtheorem{definicion}[teorema]{Definition}
\newtheorem{conjetura}[teorema]{Conjecture}
\theoremstyle{remark}
\newtheorem{observacion}[teorema]{Remark}

\newcommand{\eps}{\varepsilon}
\newcommand{\ind}{\mathbf{1}}
\DeclareMathOperator{\pdet}{pdet}
\DeclareMathOperator{\diag}{diag}
\DeclareMathOperator{\tr}{tr}
\newcommand{\Cpred}{C_{\rm pred}}
\newcommand{\Qpred}{Q_{\rm pred}}

\newcommand{\E}{\mathbb{E}}
\newcommand{\Prob}{\mathbb{P}}
\DeclareMathOperator{\per}{per}
